\section{Presupuesto de cómputo}
\label{sec:presupuesto}

\subsection{Costes unitarios medidos}

Todas las cifras salen del registro de trabajos de la campaña actual, sobre
trabajos terminados con éxito. La mediana es la cifra útil; el máximo está para que
nadie planifique con la mediana sola.

\begin{center}
\small
\begin{tabularx}{\linewidth}{@{}L{5.0cm} r r r r@{}}
\toprule
operación & $n$ & mediana & máximo & total\\
\midrule
\cod{predict\_go\_terms}       &    51 & 68.3 min & 377.8 min & 79.6 h\\
\cod{compare\_paired\_panels}  & 4{,}279 &  0.6 min &   8.3 min & 52.1 h\\
\cod{run\_cafa\_evaluation}    &   315 &  2.9 min &  48.1 min & 29.2 h\\
\cod{compute\_embeddings}      &    12 & 84.3 min & 154.6 min & 18.3 h\\
\cod{stratify\_evaluation}     &   309 &  2.9 min &  15.2 min & 18.2 h\\
\cod{load\_goa\_annotations}   &     2 & 128.8 min & 165.6 min &  4.3 h\\
\midrule
\multicolumn{4}{@{}l}{\textbf{total medido de la campaña hasta hoy}} & \textbf{202 h}\\
\bottomrule
\end{tabularx}
\captionof{table}{Coste unitario por operación. \cod{compute\_embeddings} es el
único paso que usa GPU; el KNN corre siempre en CPU.}
\end{center}

Dos observaciones sobre estas cifras antes de proyectarlas. La primera es que
\cod{predict\_go\_terms} tiene una cola larga, de 6.7 a 377.8 minutos, porque el
coste depende del tamaño del banco y de si se calculan alineamientos. La segunda es
que \cod{load\_goa\_annotations} tarda más de dos horas y es I/O contra la base:
\textbf{no se paraleliza entre máquinas}, porque una sola máquina posee el estado.

\subsection{Proyección de la rejilla cruzada}

La rejilla cruzada de sustrato por política de banco son $22 \times 16 = 352$ arms:
las 25 configuraciones menos las tres capas de \cod{ankh-base} que no son evaluables.
Con la comparación en cascada y no todos contra todos:

\begin{center}
\begin{tabularx}{\linewidth}{@{}X r r@{}}
\toprule
partida & unidades & coste\\
\midrule
codificación pendiente                   & 11 configuraciones & 120 h GPU\\
predicción                               & 352 arms          & 401 h CPU\\
evaluación, 7 cortes por 2 variantes     & 4{,}928 corridas    & 238 h CPU\\
comparación en cascada                   & unas 4{,}300        & 52 h CPU\\
\midrule
\textbf{total}                           &                   & \textbf{811 h}\\
\midrule
la misma rejilla con todos contra todos  & 61{,}776 comparaciones & +618 h, y no válida\\
\bottomrule
\end{tabularx}
\end{center}

Traducido a plazos, suponiendo cuatro trabajos concurrentes por máquina:

\begin{center}
\begin{tabularx}{\linewidth}{@{}X r r@{}}
\toprule
configuración & concurrencia & plazo\\
\midrule
un trabajador en serie            &  1 & 34 días\\
una máquina, cuatro trabajadores  &  4 & 8.5 días\\
cinco máquinas                    & 20 & \textbf{1.7 días}\\
\bottomrule
\end{tabularx}
\end{center}

\subsection{Lo que cinco máquinas desbloquean, en orden de valor}

\begin{enumerate}[leftmargin=1.5em,itemsep=0.4em]
  \item \textbf{Igualar los bancos de donantes.} Es el bloqueo que hoy impide
        cualquier veredicto de sustrato. Once configuraciones tienen 60{,}797
        embeddings frente a 528{,}294, de modo que hay que calcular unos 467{,}500
        embeddings por configuración, y son \emph{once} configuraciones. Al ritmo
        medido, entre 41{,}350 y 43{,}272 embeddings por hora, son unas 11 h de GPU
        por configuración y unas \textbf{120 h en total}, o unas \textbf{24 h sobre
        cinco GPU}. \textbf{Sigue siendo el mejor retorno del informe: desbloquea un
        eje entero por un día de cómputo sobre cinco GPU.}
  \item \textbf{La rejilla cruzada.} 1.7 días frente a 8.5. Es lo que permite
        observar interacciones entre nodos.
  \item \textbf{El encaminamiento con ajuste cruzado.} Multiplica el coste de
        puntuación por el número de pliegues. Con cinco máquinas es rutina; con una
        es el paso que se queda sin hacer, y sin él la canalización no es publicable.
  \item \textbf{La escalera de recencia del banco.} Trece releases de GOA desde la
        160 hasta la 220, de las que hay dos. \textbf{Esta no se beneficia de las
        cinco máquinas}: son unos 20.3~GB de datos y once cargas de más de dos horas
        contra la única base que posee el estado, necesariamente en serie. Son unas
        24 h de reloj y hay que planificarlas como tales.
\end{enumerate}

\begin{aviso}{Un límite que el cómputo no levanta}
Una máquina posee todo el estado: base de datos, cola y almacén de objetos. Las
otras calculan y publican de vuelta. Todo lo que sea ingesta, sellado o migración
es serie por construcción y no mejora con más nodos. Las cinco máquinas aceleran
codificación, predicción, evaluación y comparación, que son el 98\,\% del coste, y
no aceleran el 2\,\% restante. También hay una condición: los dos árboles de código
tienen que estar en la misma revisión, porque código desparejado etiqueta mal los
resultados y reporta éxito.
\end{aviso}
